% Spis skrótów (opcjonalny) - włącz przełącznikiem \showabbreviationstrue w config/metadata.tex.
% Wpisy układaj alfabetycznie. Tytuł spisu pochodzi z \abbrevname (preamble/polish.tex).
\UnnumberedChapter{\abbrevname}

\begin{abbreviations}
    \item[ACL]  Access Control List (lista kontroli dostępu)
    \item[API]  Application Programming Interface (interfejs programistyczny aplikacji)
    \item[CIDR] Classless Inter-Domain Routing (bezklasowa metoda alokacji adresów IP)
    \item[CIM]  Common Information Model (wspólny model informacji)
    \item[CLI]  Command-Line Interface (interfejs wiersza poleceń)
    \item[CSV]  Comma-Separated Values (format wartości rozdzielanych przecinkami)
    \item[GPO]  Group Policy Object (obiekt zasad grupy)
    \item[GUI]  Graphical User Interface (graficzny interfejs użytkownika)
    \item[ICMP] Internet Control Message Protocol (internetowy protokół komunikatów kontrolnych)
    \item[IP]   Internet Protocol (protokół internetowy)
    \item[JSON] JavaScript Object Notation (tekstowy format wymiany danych)
    \item[MMC]  Microsoft Management Console (konsola zarządzania Microsoft)
    \item[NLA]  Network Location Awareness (usługa rozpoznawania lokalizacji sieciowej)
    \item[OS]   Operating System (system operacyjny)
    \item[RDP]  Remote Desktop Protocol (protokół pulpitu zdalnego)
    \item[TCP]  Transmission Control Protocol (protokół kontroli transmisji)
    \item[UAC]  User Account Control (kontrola konta użytkownika)
    \item[UDP]  User Datagram Protocol (protokół pakietów użytkownika)
    \item[WFP]  Windows Filtering Platform (platforma filtrująca systemu Windows)
    \item[WMI]  Windows Management Instrumentation (instrumentacja zarządzania Windows)
    \item[XML]  Extensible Markup Language (rozszerzalny język znaczników)
\end{abbreviations}
